\section*{Discussion}

The results separate three regimes, and which one a drug falls in decides what is worth building.

Where the determinants of a drug are few and nearly sufficient, a fixed rule is as good as any
model we tried. That is the situation in \textit{Campylobacter} for all four eligible drugs, where
three well-characterised determinants account for resistance and a learned model gained nothing
($-0.006$ [$-0.017$, 0.004]), and it is close to the situation for gentamicin and cefotaxime in
external surveillance isolates, where gradient boosting lost to the rule outright. Where the
determinants are numerous, interacting or incompletely catalogued, a model on the determinant calls
adds substantially: $+0.137$ [0.097, 0.203] over the rule on unseen \Ecoli{} lineages and $+0.104$
[0.060, 0.166] externally, and a sparse linear model is enough to collect it. And where the
determinants are missing from the features altogether, as for a reference pan-genome catalogue on
unseen lineages, no amount of modelling recovers them. Prediction then leans on proxies --- lineage
membership, and assembly size as a within-lineage stand-in for accessory content --- which carry
real biology inside one collection and essentially none outside it (\nap{} 0.045 externally).

That reading explains the two-group pattern that three separate analyses on our own collection
produced: ampicillin, ciprofloxacin and trimethoprim--sulfamethoxazole lose almost nothing under a
lineage holdout once prevalence is corrected, while cefotaxime, ceftazidime,
amoxicillin--clavulanic acid and gentamicin lose 0.152--0.239 and their very major error rates rise
to between 0.68 and 0.92. Ciprofloxacin is the cleanest case of learned mechanism: a point mutation
in a target enzyme, recovered by attribution as \textit{gyrA} and \textit{parC}, learned once and
applied in an unrelated isolate. Gentamicin is the cleanest case of the opposite, and the reason it
fails has two live explanations we cannot separate. Aminoglycoside resistance is usually carried on
plasmids and other mobile elements, which travel within lineages more readily than a pan-genome
catalogue built from a fixed reference set can represent; but our own pipeline sets missing genotype
calls to the reference allele, following the source release, and such calls fail most often on
exactly those poorly assembling elements. The biological account and the imputation artefact
predict the same failure on the same drugs. Two observations keep the biological account in play:
amoxicillin--clavulanic acid, whose usual mechanism is a change in expression level rather than
gene acquisition and which therefore does not depend on assembly at all, fails in the same group;
and the drug ordering reproduces on the Giessen collection, called against a different reference.
Neither is decisive, and masking rather than imputing missing calls would settle it at the cost of
compute alone.

The pan-genome catalogue's failure is also the one place where the paper's two halves meet
directly. The drugs the catalogue loses on unseen lineages are the drugs whose determinants are
acquired, and replacing the catalogue with each strain's own determinant calls raised \nap{} on
exactly those drugs from 0.450 to 0.726 while leaving the point-mutation control unchanged. The
feature representation, not the model class, was the binding constraint. That is consistent with
work finding that gene copy-number features generalise better than SNPs under lineage shift
\cite{fistarol2025copynumber} and with the drug- and species-specificity of genotype--phenotype
discordance \cite{cheon2026discordance}.

Three consequences follow for how genomic resistance prediction should be reported, and all three
are cheap.

A determinant rule from a public catalogue should be the baseline every model is compared against,
per drug rather than on average. Our own panel contains drugs where a model beats the rule by 0.14
and drugs where the rule wins, and an average over the panel would have hidden both.

A model on metadata and assembly statistics alone should be reported on every split, because it
measures how much of an accuracy figure is available without reading a gene. On our random split
that floor was 0.705 against 0.850, and it collapsed on every split that moved populations. Without
it, the 0.850 reads as genomic skill.

Claims of generalisation need a lineage-disjoint test and, where possible, an external cohort, and
the lineage holdout must be read with the prevalence shift it causes removed. Holding out whole
clones lowers test prevalence, which lowers the metric's floor with it, so part of the resulting
drop is the measure and not the model: here two-fifths of it, which moved the gap from 0.205 to
0.118. A holdout also needs a floor on resistant strains per drug; 38
resistant strains produced an interval of [0.407, 0.524] that supports a direction and no value. The
distinction between stratifying on lineage and holding a lineage out matters here too: our
phylogroup-stratified split, which balances composition but leaves close relatives on both sides,
reported 0.856 where the true holdout reported 0.645. We built that split ourselves before
recognising it could not answer the question.

Two further points concern the model rather than the protocol. Gradient boosting transferred worse
than a sparse linear model on identical features, on six of seven external drugs, which argues for
reporting a simple learned baseline next to any complex one; nothing in the internal comparison
would have predicted it, since the two were equivalent there. And the architecture's own
interpretability did not deliver. Its gate is readable but routes ciprofloxacin away from the branch
built for it; its attention recovered no catalogued resistance gene where post-hoc attribution on
the classical learners recovered eight; and two attributions of the same pipeline share three or
four genes in twenty. We report the gene-level findings as leads for that reason, and we would not
advise treating attention weights as an explanation on a presence/absence input without the kind of
check reported here, which costs minutes and needs no retraining.

Very major errors are the finding with the most direct clinical weight, and the least comfortable.
No approach met the regulatory benchmark on the external cohort: 25.1\% for gradient boosting,
18.7\% for the linear model, 14.4\% for the rule, against an upper-limit benchmark of 7.5\%, with
amoxicillin--clavulanic acid at 35--50\% for every approach, consistent with the uncatalogued
variation reported in its resistance gene families \cite{lipworth2024uncatalogued}. The learned
models rank resistant isolates better than the rule and, at a fixed and pre-stated threshold, call
fewer of them. Moving the cut does not solve it: on our internal splits no drug met both error
limits at any threshold tested, because the cut trades one error against the other. What limits
these methods is class separation, which is to say the features, and in the drugs that matter most
the features are missing determinants rather than badly weighted.

\subsection*{Limitations}
The training collection is two-country and research-derived, with 91\% of strains from the United
Kingdom and Norway, so its metadata floor is specific to it even though the lesson is not.
Determinant calls were available for 73\% of training strains, and the pre-registered check of
whether the uncalled strains behave the same way failed (catalogue \nap{} 0.641 against 0.449), so
the rescue is established on the called subset and we do not know what the other quarter would
show. The external cohort is dominated by United States surveillance isolates and its labels come
from many laboratories; harmonising to EUCAST breakpoints helps but is not a reference-method
comparison, and our collapse of intermediate into resistant means our error rates are not the
quantities the regulatory guidance defines. Only one second species met the screening gate, and it
is one in which resistance is close to monogenic, so it bounds the advantage of a model without
telling us where that advantage ends in general. Clinical error rates were computed at a single
pre-stated threshold on uncalibrated scores.

On the model side, every architecture figure rests on a single fixed seed except the holdout margin
we repeated deliberately, so differences of a few thousandths in Tables~\ref{tab:baselines}
and~\ref{tab:attnfix} are not established effects; repeating the whole protocol over several seeds
is the obvious next step and would settle which of the smaller margins are real. The phylogroup
labels come in part from an undocumented lookup table, which is weaker than calling them from each
genome with an established in-silico typing tool \cite{beghain2018clermontyping}, and the
phylogroup-stratified split rests on them. Sequence types were taken as given, so the holdout does
not separate single-locus variants within a clonal complex and is a lower bound on lineage novelty.
Thirteen strains carry an all-zero presence/absence vector where the correct handling would extend
the label mask to that view. Missing genotype calls were imputed to the reference allele, which
biases against the mobile loci the four difficult drugs depend on. The Giessen comparison retrains
rather than transfers, and its retention figure is inflated by that collection's higher prevalence.
None of the learned models compared here was tuned; better ones exist, and the comparisons are fair
only in the narrow sense that every model saw the same features and folds.
